\documentclass[10pt,aspectratio=169]{beamer}

% ---------- Theme ----------
\usetheme{Madrid}
\usecolortheme{seahorse}
\setbeamertemplate{navigation symbols}{}
\setbeamertemplate{footline}[frame number]

% ---------- Packages ----------
\usepackage[utf8]{inputenc}
\usepackage{amsmath,amssymb,amsthm}
\usepackage{braket}
\usepackage{bm}
\usepackage{booktabs}
\usepackage{listings}
\usepackage{xcolor}
\usepackage{hyperref}

% ---------- Code listing style ----------
\definecolor{codebg}{rgb}{0.96,0.96,0.98}
\definecolor{codekw}{rgb}{0.10,0.30,0.70}
\definecolor{codestr}{rgb}{0.55,0.20,0.20}
\definecolor{codecom}{rgb}{0.25,0.55,0.25}

\lstdefinestyle{pystyle}{
  language=Python,
  basicstyle=\ttfamily\scriptsize,
  keywordstyle=\color{codekw}\bfseries,
  stringstyle=\color{codestr},
  commentstyle=\color{codecom}\itshape,
  backgroundcolor=\color{codebg},
  frame=single,
  rulecolor=\color{black!20},
  showstringspaces=false,
  breaklines=true,
  numbers=left,
  numberstyle=\tiny\color{gray},
  tabsize=2,
  columns=fullflexible,
  upquote=true,
}
\lstset{style=pystyle}

% ---------- Title info ----------
\title[VQE-UCCSD from a Hartree--Fock reference]
{From a Hartree--Fock reference to a VQE-UCCSD calculation}
\subtitle{A hybrid quantum--classical workflow for electronic ground states}
\author{Morten Hjorth-Jensen}
\date{\today}

\begin{document}

% =========================================================
\begin{frame}
\titlepage
\end{frame}

% =========================================================
\begin{frame}{Outline}
\tableofcontents
\end{frame}

% =========================================================
\section{The eigenvalue problem and the VQE strategy}

\begin{frame}{Goal: the electronic ground-state energy}
We want to approximate the lowest eigenvalue of an electronic Hamiltonian,
\[
E_0 \;=\; \min_{\ket{\Psi}}\;
\frac{\bra{\Psi}\hat H\ket{\Psi}}{\braket{\Psi|\Psi}}.
\]

Exact diagonalisation scales as $\mathcal{O}(d^3)$ with Hilbert-space dimension $d=2^{N_q}$ --- intractable for moderate molecules.

\medskip
The \textbf{Variational Quantum Eigensolver (VQE)} restricts the wave function to a parameterised ansatz,
\[
\ket{\Psi(\bm\theta)} \;=\; U(\bm\theta)\,\ket{\Phi_0},
\]
and minimises
\[
E(\bm\theta) \;=\; \bra{\Psi(\bm\theta)}\hat H\ket{\Psi(\bm\theta)} \;\geq\; E_0.
\]

For UCCSD the reference state is the Hartree--Fock determinant,
$\ket{\Phi_0}=\ket{\Phi_{\mathrm{HF}}}$. The rest of these notes
develop \emph{why} this is a good choice and \emph{how} to realise it on a quantum computer.
\end{frame}

% =========================================================
\begin{frame}{Why the initial state matters}
The ansatz is $\ket{\Psi(\bm\theta)} = U(\bm\theta)\ket{\Phi_0}$. The reference state $\ket{\Phi_0}$ controls:
\begin{itemize}
  \item the region of Hilbert space reachable at a given circuit depth,
  \item the number of variational parameters needed for chemical accuracy,
  \item susceptibility to barren plateaus and local minima.
\end{itemize}

\medskip
A bad reference (e.g.\ the qubit vacuum $\ket{0\cdots 0}$) forces $U(\bm\theta)$ first to build in the gross electronic structure before refining correlation.

\medskip
A good reference encodes the \emph{mean-field} physics already, so $U(\bm\theta)$ only has to add electron correlation on top.

\medskip
$\Rightarrow$ the Hartree--Fock state is the natural starting point inherited from classical quantum chemistry.
\end{frame}

% =========================================================
\section{The electronic Hamiltonian}

\begin{frame}{The electronic Hamiltonian in second quantisation}
In second quantisation,
\[
\hat H \;=\; \sum_{pq} h_{pq}\, \hat a_p^\dagger \hat a_q
\;+\;\frac12\sum_{pqrs} h_{pqrs}\, \hat a_p^\dagger \hat a_q^\dagger \hat a_s \hat a_r,
\]
with one-body integrals
\[
h_{pq} \;=\; \int d\mathbf{x}\,
\phi_p^*(\mathbf{x})\!\left(-\tfrac12\nabla^2 - \sum_A \tfrac{Z_A}{|\mathbf{r}-\mathbf{R}_A|}\right)\!\phi_q(\mathbf{x}),
\]
and two-body integrals
\[
h_{pqrs} \;=\;
\int d\mathbf{x}_1\,d\mathbf{x}_2\;
\frac{\phi_p^*(\mathbf{x}_1)\phi_q^*(\mathbf{x}_2)\phi_r(\mathbf{x}_1)\phi_s(\mathbf{x}_2)}{|\mathbf{r}_1-\mathbf{r}_2|}.
\]

\smallskip
The single-particle basis $\{\phi_p\}$ is typically obtained from a preceding Hartree--Fock calculation,
which is the next ingredient.
\end{frame}

% =========================================================
\section{The Hartree--Fock reference}

\begin{frame}{Hartree--Fock recap}
The Hartree--Fock (HF) method approximates the many-electron wave function by a single Slater determinant,
\[
\ket{\Phi_{\mathrm{HF}}} \;=\; \frac{1}{\sqrt{N_e!}}\,
\det\!\big[\,\varphi_{i_1}(x_1)\cdots\varphi_{i_{N_e}}(x_{N_e})\,\big],
\]
with spin-orbitals $\{\varphi_p\}$ obtained self-consistently from the Roothaan--Hall equations.

\medskip
In second quantisation this is just a product of creation operators acting on the Fock vacuum,
\[
\boxed{\;
\ket{\Phi_{\mathrm{HF}}} \;=\; \prod_{p\,\in\,\mathrm{occ}} \hat a_p^\dagger \,\ket{\mathrm{vac}}\;}
\]
where $\mathrm{occ}$ contains the $N_e$ lowest-energy spin-orbitals.

\medskip
HF captures the mean-field physics exactly. What it misses is the \emph{correlation energy},
\[
E_{\mathrm{corr}} \;=\; E_{\mathrm{exact}} - E_{\mathrm{HF}}.
\]
Recovering $E_{\mathrm{corr}}$ on a quantum computer is precisely the job of $U(\bm\theta)$.
\end{frame}

% =========================================================
\section{Mapping fermions to qubits}

\begin{frame}{From fermions to qubits}
A qubit can encode the occupation of one spin-orbital,
\[
\ket{0}_p \leftrightarrow \text{empty},\qquad \ket{1}_p \leftrightarrow \text{occupied}.
\]
We need a mapping that respects fermionic antisymmetry. The fermionic Hamiltonian becomes
\[
\hat H \;\longrightarrow\; \hat H_q \;=\; \sum_\alpha c_\alpha\, P_\alpha,
\qquad
P_\alpha = \sigma_1^{(\alpha)}\otimes\sigma_2^{(\alpha)}\otimes\cdots\otimes\sigma_n^{(\alpha)},
\]
a weighted sum of Pauli strings. Common mappings:
\begin{itemize}
  \item \textbf{Jordan--Wigner (JW):} simple, local in occupation, non-local parity strings.
  \item \textbf{Parity:} swaps the locality trade-off, allows a two-qubit reduction via $\mathbb{Z}_2$ symmetries.
  \item \textbf{Bravyi--Kitaev (BK):} $\mathcal{O}(\log N)$ locality in both occupation and parity.
\end{itemize}

\smallskip
Under Jordan--Wigner,
\[
\hat a_p^\dagger \;=\; \Big(\prod_{q<p} Z_q\Big)\,\frac{X_p - i Y_p}{2}.
\]

\smallskip
The VQE objective then reads
\(\,E(\bm\theta) = \sum_\alpha c_\alpha \langle P_\alpha\rangle_{\bm\theta}.\)
\end{frame}

% =========================================================
\begin{frame}{The HF state under Jordan--Wigner}
With JW ordering (spin-orbitals indexed $0,1,\ldots,M-1$, occupied first),
\[
\ket{\Phi_{\mathrm{HF}}}_{\mathrm{JW}}
\;=\;
\underbrace{\ket{1}\ket{1}\cdots\ket{1}}_{N_e\ \text{occupied}}\,
\underbrace{\ket{0}\ket{0}\cdots\ket{0}}_{M-N_e\ \text{virtual}}.
\]

\textbf{This is a product state!} It is prepared from $\ket{0}^{\otimes M}$ by applying $X$ gates on the first $N_e$ qubits.

\medskip
For $\mathrm{H}_2$ in the minimal STO-3G basis,
\[
M=4\text{ spin-orbitals},\quad N_e=2\;\;\Longrightarrow\;\;\ket{\Phi_{\mathrm{HF}}}=\ket{1100}.
\]

\medskip
Different mappings produce different bit patterns for the \emph{same} physical state:
\begin{center}
\begin{tabular}{lll}
\toprule
Mapping & H$_2$ HF state & \# qubits \\
\midrule
Jordan--Wigner            & $\ket{1100}$ & 4 \\
Parity                    & $\ket{0110}$ & 4 \\
Parity $+\,\mathbb{Z}_2$  & $\ket{01}$   & 2 \\
Bravyi--Kitaev            & $\ket{1110}$ & 4 \\
\bottomrule
\end{tabular}
\end{center}
Never hard-code the bit pattern --- derive it from the mapper.
\end{frame}

% =========================================================
\section{Preparing the HF state in code}

\begin{frame}[fragile]{A bare-bones HF initialiser in Qiskit}
Under Jordan--Wigner, preparing $\ket{\Phi_{\mathrm{HF}}}$ takes $N_e$ Pauli-$X$ gates:
\begin{lstlisting}
from qiskit import QuantumCircuit, QuantumRegister

def hartree_fock_jw(n_spin_orbitals, n_electrons):
    """Prepare a Hartree-Fock reference state under JW mapping.
    Occupied orbitals are the lowest n_electrons indices."""
    qr = QuantumRegister(n_spin_orbitals, name="q")
    qc = QuantumCircuit(qr, name="HF")
    for i in range(n_electrons):
        qc.x(qr[i])           # flip |0> -> |1> on occupied orbitals
    return qc

hf_circuit = hartree_fock_jw(n_spin_orbitals=4, n_electrons=2)
print(hf_circuit)
\end{lstlisting}

\smallskip
Schematic output:
\begin{verbatim}
q_0: -X-
q_1: -X-
q_2: ---
q_3: ---
\end{verbatim}
That is the entire HF state preparation under JW.
\end{frame}

% =========================================================
\begin{frame}[fragile]{Using Qiskit Nature's built-in HF initial state}
For production work, let the library handle JW, parity (with two-qubit reduction), and BK consistently:
\begin{lstlisting}
from qiskit_nature.second_q.circuit.library import HartreeFock
from qiskit_nature.second_q.mappers import JordanWignerMapper

mapper = JordanWignerMapper()

hf_init = HartreeFock(
    num_spatial_orbitals = 2,       # H2 in STO-3G has 2 spatial orbitals
    num_particles        = (1, 1),  # (alpha, beta) electrons
    qubit_mapper         = mapper,
)
print(hf_init.draw())
\end{lstlisting}

\smallskip
Switching mapper $=$ switching the bit-pattern that represents HF; the library does the bookkeeping.
With $\ket{\Phi_{\mathrm{HF}}}$ in hand, we can now ask the quantum computer to add the correlation.
\end{frame}

% =========================================================
\section{From Hartree--Fock to coupled cluster}

\begin{frame}{Classical coupled cluster}
In classical coupled-cluster (CC) theory, the correlated wave function is
\[
\ket{\Psi_{\mathrm{CC}}} \;=\; e^{\hat T}\,\ket{\Phi_{\mathrm{HF}}},
\qquad
\hat T \;=\; \hat T_1 + \hat T_2 + \hat T_3 + \cdots
\]
with single and double excitation operators
\[
\hat T_1 \;=\; \sum_{ia} t_i^{a}\,\hat a_a^\dagger \hat a_i,
\qquad
\hat T_2 \;=\; \frac14\sum_{ijab} t_{ij}^{ab}\,\hat a_a^\dagger \hat a_b^\dagger \hat a_j \hat a_i.
\]
Indices: $i,j$ are occupied spin-orbitals; $a,b$ are virtual spin-orbitals.

\bigskip
Each excitation promotes one (or two) electrons from below to above the Fermi level of the HF reference.
The coefficients $t_i^a,\,t_{ij}^{ab}$ are the cluster amplitudes that encode electron correlation.

\bigskip
\textbf{Problem for quantum hardware:} the operator $e^{\hat T}$ is \emph{not} unitary, so it cannot be implemented directly as a quantum circuit.
\end{frame}

% =========================================================
\begin{frame}{Unitary coupled cluster}
The fix is to use an anti-Hermitian generator. Define
\[
\hat A \;=\; \hat T - \hat T^\dagger,
\qquad \hat A^\dagger = -\hat A.
\]
Then $e^{\hat A}$ \emph{is} unitary, and the unitary coupled-cluster ansatz reads
\[
\ket{\Psi_{\mathrm{UCC}}}
\;=\; e^{\hat T - \hat T^\dagger}\,\ket{\Phi_{\mathrm{HF}}}.
\]

Truncating at singles and doubles gives \textbf{UCCSD},
\[
\ket{\Psi_{\mathrm{UCCSD}}}
\;=\; e^{\hat T_1 + \hat T_2 - \hat T_1^\dagger - \hat T_2^\dagger}\,\ket{\Phi_{\mathrm{HF}}}.
\]

\bigskip
Two crucial properties that justify starting from the HF reference:
\begin{itemize}
  \item At $\bm\theta=0$ all amplitudes vanish and $\ket{\Psi(0)}=\ket{\Phi_{\mathrm{HF}}}$.
        The optimiser starts at the (known) HF energy --- a sensible baseline.
  \item Each excitation operator promotes electrons from occupied to virtual orbitals,
        i.e.\ it generates correlation \emph{on top of} HF.
\end{itemize}
\end{frame}

% =========================================================
\begin{frame}{Single and double excitation operators}
\textbf{Single excitation} $i\to a$ has anti-Hermitian generator
\[
\hat\tau_i^{a} \;=\; \hat a_a^\dagger \hat a_i \;-\; \hat a_i^\dagger \hat a_a,
\]
giving the unitary
\[
U_i^{a}(\theta_i^a) \;=\; \exp\!\big[\,\theta_i^a\big(\hat a_a^\dagger \hat a_i - \hat a_i^\dagger \hat a_a\big)\big].
\]
This rotates between the reference and singly excited determinants.

\bigskip
\textbf{Double excitation} $ij\to ab$ has generator
\[
\hat\tau_{ij}^{ab}
\;=\;
\hat a_a^\dagger \hat a_b^\dagger \hat a_j \hat a_i
\;-\;
\hat a_i^\dagger \hat a_j^\dagger \hat a_b \hat a_a,
\]
and unitary
\[
U_{ij}^{ab}(\theta_{ij}^{ab})
\;=\;
\exp\!\big[\,\theta_{ij}^{ab}\,\hat\tau_{ij}^{ab}\big].
\]
Doubles are essential for capturing dynamical electron correlation.
\end{frame}

% =========================================================
\begin{frame}{Trotterised UCCSD ansatz}
The exact UCCSD unitary is
\[
U_{\mathrm{UCCSD}}(\bm\theta)
\;=\;
\exp\!\Big[\,\sum_{ia}\theta_i^a\,\hat\tau_i^a \;+\; \sum_{ijab}\theta_{ij}^{ab}\,\hat\tau_{ij}^{ab}\Big].
\]

Excitation generators generally do not commute, $[\hat\tau_\mu,\hat\tau_\nu]\neq 0$,
so a single exponential cannot be implemented exactly. In practice one uses a first-order Trotter approximation,
\[
U_{\mathrm{UCCSD}}(\bm\theta)
\;\approx\;
\prod_{ia} e^{\theta_i^a\,\hat\tau_i^a}\;
\prod_{ijab} e^{\theta_{ij}^{ab}\,\hat\tau_{ij}^{ab}}.
\]
Each factor is a Pauli-string rotation that compiles to a short sequence of CNOT and single-qubit gates.

\bigskip
The Trotter error is absorbed by the variational optimisation:
the parameters $\bm\theta$ are re-tuned to minimise the energy of whatever circuit is actually executed.
\end{frame}

% =========================================================
\section{Building and running the VQE-UCCSD circuit}

\begin{frame}{Quantum circuit construction}
The full UCCSD trial state on a quantum computer is
\[
\ket{\Psi(\bm\theta)}
\;=\;
U_{\mathrm{UCCSD}}(\bm\theta)\,\ket{\Phi_{\mathrm{HF}}}.
\]

\textbf{Workflow:}
\begin{enumerate}
  \item Prepare $\ket{\Phi_{\mathrm{HF}}}$ on the qubits ($N_e$ Pauli-$X$ gates under JW).
  \item Apply all selected single-excitation unitaries $\prod_{ia} e^{\theta_i^a\hat\tau_i^a}$.
  \item Apply all selected double-excitation unitaries $\prod_{ijab} e^{\theta_{ij}^{ab}\hat\tau_{ij}^{ab}}$.
  \item Measure expectation values of the Pauli strings in $\hat H_q$.
\end{enumerate}

\medskip
The HF preparation in step 1 is essentially free (a single layer of $X$ gates).
Steps 2--3 dominate the depth: UCCSD circuits contain many excitations and become deep for larger systems.
\end{frame}

% =========================================================
\begin{frame}{Energy measurement}
After fermion-to-qubit mapping the Hamiltonian is
\[
\hat H_q \;=\; \sum_\alpha c_\alpha\, P_\alpha.
\]
For a fixed parameter vector $\bm\theta$ we estimate
\[
E(\bm\theta)
\;=\;
\sum_\alpha c_\alpha \bra{\Psi(\bm\theta)} P_\alpha \ket{\Psi(\bm\theta)}.
\]

\medskip
Each Pauli string is measured separately on the quantum device. For example,
\[
P_\alpha \;=\; Z_0\, Z_1\, X_2\, Y_3,
\]
requires rotating qubits 2 and 3 into the $X$ and $Y$ bases (Hadamard, $S^\dagger H$) before a $Z$-basis measurement.

\medskip
The number of distinct Pauli strings grows polynomially with system size, and grouping commuting terms reduces the number of measurement circuits.
\end{frame}

% =========================================================
\begin{frame}{The hybrid quantum--classical loop}
VQE is hybrid: each iteration combines quantum estimation with classical optimisation.
\[
\bm\theta^{(k)}
\;\longrightarrow\;
\text{quantum circuit}
\;\longrightarrow\;
E\big(\bm\theta^{(k)}\big)
\;\longrightarrow\;
\text{classical optimiser}.
\]

The optimiser updates the parameters,
\[
\bm\theta^{(k+1)} \;=\; \bm\theta^{(k)} + \Delta\bm\theta^{(k)},
\]
until convergence,
\[
\big|E^{(k+1)} - E^{(k)}\big| \;<\; \epsilon.
\]

\medskip
The variational principle guarantees
\[
E_{\mathrm{VQE}} \;=\; \min_{\bm\theta} E(\bm\theta) \;\geq\; E_0,
\]
so improving the ansatz (more excitations, deeper Trotter, larger basis) can only lower the variational energy.
\end{frame}

% =========================================================
\begin{frame}{Gradients and the parameter-shift rule}
Gradient-free optimisers (COBYLA, SPSA, SLSQP) are common, but gradients can also be computed on hardware.
For a parameter $\theta_\mu$,
\[
\frac{\partial E}{\partial \theta_\mu}
\;=\;
\bra{\partial_\mu \Psi}\hat H\ket{\Psi} + \bra{\Psi}\hat H\ket{\partial_\mu \Psi}.
\]

For Pauli-generated rotations
\(\,U_\mu(\theta_\mu) = e^{-i\theta_\mu P_\mu/2},\)
the \textbf{parameter-shift rule} gives an exact analytic gradient from two extra circuit evaluations:
\[
\frac{\partial E}{\partial \theta_\mu}
\;=\;
\frac12\left[\,E\!\Big(\theta_\mu+\tfrac{\pi}{2}\Big) - E\!\Big(\theta_\mu-\tfrac{\pi}{2}\Big)\right].
\]

\medskip
This is structurally similar to a finite difference but is \emph{exact} for these gate generators --- no step-size bias.
\end{frame}

% =========================================================
\section{The complete algorithm}

\begin{frame}{Algorithm: VQE-UCCSD with a Hartree--Fock reference}
\small
\begin{enumerate}
  \item Choose a molecular geometry and a one-particle basis (e.g.\ STO-3G).
  \item Compute one- and two-body integrals $h_{pq}$ and $h_{pqrs}$ (PySCF, Psi4, ...).
  \item Assemble the second-quantised Hamiltonian
        \[
        \hat H = \sum_{pq} h_{pq}\hat a_p^\dagger \hat a_q + \tfrac12 \sum_{pqrs} h_{pqrs}\hat a_p^\dagger \hat a_q^\dagger \hat a_s \hat a_r.
        \]
  \item Map to a qubit Hamiltonian via JW / parity / BK:
        $\hat H_q = \sum_\alpha c_\alpha P_\alpha$.
  \item Prepare the Hartree--Fock reference $\ket{\Phi_{\mathrm{HF}}}$.
  \item Build the UCCSD ansatz on top of it:
        \(U(\bm\theta) = e^{\hat T_1 + \hat T_2 - \hat T_1^\dagger - \hat T_2^\dagger}.\)
  \item Initialise $\bm\theta = 0$ (so the trial state starts at $\ket{\Phi_{\mathrm{HF}}}$).
  \item Measure $E(\bm\theta) = \bra{\Psi(\bm\theta)}\hat H_q\ket{\Psi(\bm\theta)}$.
  \item Update $\bm\theta$ with a classical optimiser; repeat 8--9 until $|\Delta E| < \epsilon$.
\end{enumerate}
\end{frame}

% =========================================================
\begin{frame}[fragile]{Putting it all together: H$_2$ with Qiskit Nature}
\begin{lstlisting}
from qiskit_nature.second_q.drivers import PySCFDriver
from qiskit_nature.second_q.mappers import JordanWignerMapper
from qiskit_nature.second_q.circuit.library import HartreeFock, UCCSD
from qiskit_algorithms           import VQE
from qiskit_algorithms.optimizers import SLSQP
from qiskit.primitives           import Estimator
import numpy as np

# 1-3. Electronic structure problem and second-quantised Hamiltonian
driver  = PySCFDriver(atom="H 0 0 0; H 0 0 0.735", basis="sto3g")
problem = driver.run()

# 4. Map fermions -> qubits
mapper   = JordanWignerMapper()
qubit_op = mapper.map(problem.hamiltonian.second_q_op())

# 5. Hartree-Fock initial state
hf_state = HartreeFock(problem.num_spatial_orbitals,
                       problem.num_particles, mapper)

# 6. UCCSD ansatz built ON TOP OF the HF reference
ansatz = UCCSD(problem.num_spatial_orbitals,
               problem.num_particles, mapper,
               initial_state=hf_state)

# 7-9. Run VQE starting from theta = 0  (state = |Phi_HF>)
vqe = VQE(Estimator(), ansatz, SLSQP(maxiter=200),
          initial_point=np.zeros(ansatz.num_parameters))
result = vqe.compute_minimum_eigenvalue(qubit_op)

print("VQE electronic energy:", result.eigenvalue.real)
\end{lstlisting}
\end{frame}

% =========================================================
\begin{frame}{What this code is doing, line by line}
\begin{itemize}
  \item \textbf{Driver + problem:} PySCF computes the HF orbitals and the integrals $h_{pq}$, $h_{pqrs}$ in that basis.
  \item \textbf{Mapper:} converts $\hat H$ into a sum of Pauli strings $\hat H_q$.
  \item \textbf{HartreeFock:} a circuit that prepares $\ket{\Phi_{\mathrm{HF}}}$ --- under JW just $N_e$ Pauli-$X$ gates.
  \item \textbf{UCCSD(initial\_state=hf\_state):} the trial state is $U_{\mathrm{UCCSD}}(\bm\theta)\ket{\Phi_{\mathrm{HF}}}$.
  \item \textbf{initial\_point=zeros:} at iteration $0$, $\bm\theta=0$, so $\ket{\Psi(0)} = \ket{\Phi_{\mathrm{HF}}}$ and $E(0) = E_{\mathrm{HF}}$.
  \item \textbf{VQE:} the hybrid loop runs --- the optimiser only has to find the \emph{correlation} energy
        $E_{\mathrm{corr}} = E_{\mathrm{exact}} - E_{\mathrm{HF}}$, a small fraction of the total.
\end{itemize}

\medskip
For H$_2$/STO-3G this converges in a handful of iterations and reproduces the FCI energy to machine precision.
\end{frame}

% =========================================================
\section{Properties, scaling, and limitations}

\begin{frame}{Scaling of UCCSD}
Let
\[
N_{\mathrm{occ}} = \text{\# occupied spin-orbitals},
\qquad
N_{\mathrm{virt}} = \text{\# virtual spin-orbitals}.
\]

Number of excitation amplitudes:
\[
N_{\mathrm{singles}} \;\sim\; N_{\mathrm{occ}}\,N_{\mathrm{virt}},
\qquad
N_{\mathrm{doubles}} \;\sim\; N_{\mathrm{occ}}^2\,N_{\mathrm{virt}}^2.
\]

\medskip
Each amplitude contributes a Trotter factor $e^{\theta_\mu \hat\tau_\mu}$, which compiles to multiple CNOTs.
The circuit depth therefore grows roughly as $\mathcal{O}(N_{\mathrm{occ}}^2 N_{\mathrm{virt}}^2)$,
which becomes prohibitive for medium-sized molecules on near-term devices.

\medskip
Mitigations: active-space selection, point-group/spin symmetries, k-UpCCGSD, ADAPT-VQE (grow the operator pool adaptively), tapering off $\mathbb{Z}_2$ symmetries (e.g.\ two-qubit reduction).
\end{frame}

% =========================================================
\begin{frame}{Why UCCSD is useful --- and where it hurts}
\textbf{Strengths:}
\begin{itemize}
  \item Physically motivated; direct link to classical quantum chemistry.
  \item Captures dynamical electron correlation systematically.
  \item Unitary form is natural for quantum circuits.
  \item Size-extensive in the untruncated limit.
\end{itemize}

\medskip
\textbf{Limitations:}
\begin{itemize}
  \item Deep circuits and many variational parameters.
  \item Trotter errors (mitigated by re-optimising $\bm\theta$).
  \item Heavy measurement overhead --- thousands of Pauli strings for larger molecules.
  \item Optimisation landscape can suffer from local minima and barren plateaus.
  \item HF is a poor reference for strongly correlated systems (stretched bonds, transition metals).
\end{itemize}
\end{frame}

% =========================================================
\begin{frame}{Beyond a single-determinant reference}
When HF fails (strong static correlation), the same machinery still applies but with a richer $\ket{\Phi_0}$:
\begin{itemize}
  \item \textbf{Multi-determinant references} (CASSCF, CISD): prepared by an entangling layer of controlled rotations on top of the HF X-gates.
  \item \textbf{Symmetry-broken (UHF) orbitals}: cheaper, but lose spin symmetry.
  \item \textbf{ADAPT-VQE}: grow the operator pool one excitation at a time, partly compensating for a poor reference.
\end{itemize}

\medskip
The structure of the VQE-UCCSD workflow does not change --- only the first circuit block (initial state) and possibly the operator pool. Everything else (mapping, measurement, optimisation, parameter-shift gradients) carries over unchanged.
\end{frame}

% =========================================================
\section{Summary}

\begin{frame}{Conceptual summary}
The VQE-UCCSD method combines
\[
\text{quantum chemistry}\;+\;\text{unitary coupled cluster}\;+\;\text{quantum measurement}\;+\;\text{classical optimisation}.
\]

The central variational problem is
\[
E_{\mathrm{VQE}}
\;=\;
\min_{\bm\theta}\,
\bra{\Phi_{\mathrm{HF}}}
U^\dagger_{\mathrm{UCCSD}}(\bm\theta)\,
\hat H\,
U_{\mathrm{UCCSD}}(\bm\theta)
\ket{\Phi_{\mathrm{HF}}}.
\]

\medskip
\textbf{Take-home messages:}
\begin{itemize}
  \item The Hartree--Fock state is the natural reference; under Jordan--Wigner it is a single product state $\ket{1\cdots1\,0\cdots0}$.
  \item Starting VQE at $\bm\theta=0$ means $\ket{\Psi(0)} = \ket{\Phi_{\mathrm{HF}}}$ and $E(0) = E_{\mathrm{HF}}$; only the correlation energy needs to be recovered.
  \item UCCSD adds correlation through anti-Hermitian single and double excitations, $e^{\hat T - \hat T^\dagger}$, which is unitary and circuit-implementable.
  \item The quantum computer prepares the correlated wave function; the classical computer updates the cluster amplitudes.
  \item For strongly correlated systems, replace HF with a multi-determinant or adaptive reference --- the rest of the pipeline is unchanged.
\end{itemize}
\end{frame}

% =========================================================
\begin{frame}{References}
\footnotesize
\begin{thebibliography}{99}
\bibitem{Peruzzo2014}
A. Peruzzo \emph{et al.},
A variational eigenvalue solver on a photonic quantum processor,
\emph{Nat.\ Commun.} \textbf{5}, 4213 (2014).

\bibitem{McClean2016}
J. R. McClean \emph{et al.},
The theory of variational hybrid quantum-classical algorithms,
\emph{New J.\ Phys.} \textbf{18}, 023023 (2016).

\bibitem{Romero2018}
J. Romero \emph{et al.},
Strategies for quantum computing molecular energies using the unitary coupled cluster ansatz,
\emph{Quantum Sci.\ Technol.} \textbf{4}, 014008 (2018).

\bibitem{Bartlett2007}
R. J. Bartlett and M. Musial,
Coupled-cluster theory in quantum chemistry,
\emph{Rev.\ Mod.\ Phys.} \textbf{79}, 291 (2007).

\bibitem{McArdle2020}
S. McArdle \emph{et al.},
Quantum computational chemistry,
\emph{Rev.\ Mod.\ Phys.} \textbf{92}, 015003 (2020).

\bibitem{Tilly2022}
J. Tilly \emph{et al.},
The Variational Quantum Eigensolver: a review of methods and best practices,
\emph{Phys.\ Rep.} \textbf{986}, 1 (2022).
\end{thebibliography}
\end{frame}

\end{document}
